There is a large and growing literature in programming frameworks for LLMs, foundation models, and other AI artifacts.

Many current frameworks are focused on prompt management. 
LangChain~\cite{langchain} and LlamaIndex~\cite{Liu_LlamaIndex_2022} are popular libraries for building LLM applications, with special support for mechanical issues surrounding AI model use: managing prompt templates, handling user-provided labeled examples, managing document tokenization, and so on. They have no higher-level task representation, although there is support for standard use cases like chatbots and RAG.

DSPy~\cite{dspy} focuses on creating high-quality prompts. It takes a high-level description of a user's LLM goal and yields a set of concrete prompts (or even LLM operations like fine-tuning) to implement the user's goal. Like \system, DSPy aims to abstract some details of the AI programming process. But DSPy's goals are different: it focuses entirely on improving prompt quality, ignores most performance and cost issues, and asks the programmer to make decisions about its machine learning algorithm.  We have used DSPy in some cases when compiling individual tasks inside \system.

Several other projects --- such as Outlines~\cite{willard2023efficient}, Guidance~\cite{guidance_ai_2023}, and RELM~\cite{kuchnik2023validating} --- attempt to control how LLMs emit data. These projects are especially useful when producing data that is intended for machine consumption, such as producing a JSON record that conforms to a particular schema, or emitting classification labels drawn from a target vocabulary. Outlines and RELM offer runtime optimizations that avoid inference costs when the syntax fully determines the output. These are useful systems but are also fairly narrow and do not address our target workload. 

SGLang~\cite{zheng2023efficiently} is an ambitious project that combines a number of features present in the above systems, such as prompt templating methods and constrained outputs. It also tries to offer some performance features not present in the above systems, such as parallel execution and batching. It even offers management of multimodal data. However, the Structured Generation Language is still fairly low-level and asks developers to make many decisions manually, such as the direct text of a prompt, the desired elements that comprise a batch of  prompts; or the degree of parallelism. It might make sense for a future version of \system\ to use SGLang as a compiled runtime description language.

The SkyPilot system~\cite{Yang2023SkyPilotAI} is similar to \system\ in its emphasis on cost reduction of ML workloads, but is primarily concerned with efficiently deploying coarse DAG of tasks across broadly-similar cloud platforms.  It offers a broker that stands between a user's workload and advertised services from competing cloud services. It aims for an ideal assignment of work to resources, rather than formulating efficient program-specific implementations.

FrugalGPT~\cite{chen2023frugalgpt} is a high-level library meant to be deployed on top of LLM access packages such as OpenAI~\cite{openaiapi}. It offers a number of valuable optimizations, some of which resemble the optimizations in \autoref{sec:optimizations}. The FrugalGPT vision of appealing AI optimizations, combined with some trade-offs, is entirely consistent with \system's goal and approach, and the observed improvements in runtime and cost are strong. However, it is unclear how the user describes tasks to the FrugalGPT's optimizer, so the scope for future optimization is also unclear. Moreover, it is not clear whether FrugalGPT optimizations can be implemented in the context of a broader AI program that includes not just LLM processing but also non-LLM AI components and conventional data processing elements.

AutoGen~\cite{wu2023autogen} is an interesting development framework for "LLM applications," but with a domain focus that is dramatically different from the one we pursue. This system imagines LLM applications as conversations among agents, which might be humans, LLMs, or tools. The AutoGen developer aims to design the conversation and control flow that takes place among these agents. This vision of LLM applications is driven by a fundamentally different model from our own, which is based more in the data processing tradition.

Urban and Binnig proposed "Language-Model-Driven Query Planning" in the Caesura system~\cite{urban2023caesura}. The goal is to allow users to describe a data processing task in natural language, which yields an executable query plan, in particular one that can operate on multimodal data. This work is similar to \system\ in that it yields a query plan that combines traditional relational operators with model-driven ones. However, \system\ makes no attempt at all to build the logical query plan from natural language; we expect a human programmer to provide the core processing logic. The published Caesura system does not perform any plan optimization, although the authors describe this as an area for future work.

ZenDB \cite{lin2024towards} is a document analytics system that enhances query performance by extracting semantic hierarchical structures from templatized documents and integrating a query engine that leverages these structures to efficiently and accurately execute SQL queries on document collections. ZenDB employs optimization techniques such as predicate reordering, pushdown, and projection pull-up to refine the execution of user-specified SQL queries. This optimization is facilitated through a summary-based tree search for each document, aiming to minimize cost and latency while maximizing accuracy. Although ZenDB and similar systems share optimization objectives with \system, their efforts primarily concentrate on logical optimization and are predominantly focused on managing text data. In contrast, \system\ utilizes multiple strategies to optimize broader aspects of query execution.

The Evaporate system proposed by Arora et al.\ in~\cite{arora2023language} shares some intuitions of \system{} regarding the trade-off of cost/quality in LLM workloads using code generation.
However, while Evaporate is designed specifically for an information extraction use case and employs only static prompts and code generation, \system{} provides a more general approach to express a wider class of ML workloads, with several optimization strategies alongside code generation, e.g., model selection and token reduction.
As such, we believe the more advanced weak-supervision code generation strategy proposed in~\cite{arora2023language} could be integrated as one of the optimizations within the \system{} framework.  


%For example, if the program involves textual filtering, \system may choose to employ a Regex expression instead of calling a LLM.
%\system can integrate different optimization modules, in facts some of its %concrete operators are currently sometimes compiled into DSPy programs. 

%Other systems handle parts of the neural model data processing framework, but do not focus heavily on prompt mangement. LlamaIndex~\cite{llamaindex} is a python framework for easily constructing context for Retrieval Augmented Generation tasks. A number of vector databases --- which offer fast querying of embeddings, and thus might power a RAG system --- exist. RAG components might be marshalled as part of a \system\ program. But these systems do not attempt to express full AI programs or perform declarative query optimization.  (REMIND: this description of LLamaIndex might be out of date, need to review)



\system{} has some commonalities with  CrowdDB~\cite{crowddb}, Qurk~\cite{qurk}, and Deco~\cite{deco};  systems that aim to provide a declarative approach to create data processing pipelines with crowd-sourcing operations to answer queries that cannot ``be answered by machines only'' \cite{crowddb}. 
For example, ~\cite{crowddb} automatically optimized data processing pipelines with  crowd-sourcing operations for entity matching or (open) data retrieval, whereas ~\cite{crowdcount} introduced and optimized crowd-sourcing operators to extract data from images. 
It turns out that --- with the power of LLMs --- some of the tasks no longer need humans and can be answered by machines only.
Instead of using crowd-sourcing operators which give humans tasks to do, LLMs can perform tasks like extracting data from text, images, etc. much faster and often with better quality. 
However, any task done by an LLM still has a quality/time/cost trade-off and, similar to crowd-sourcing tasks, the quality increases the more one is willing to pay per task (e.g., by using a large vs small model, single vs multiple invocations, etc.). 
However, while this trade-off creates some similarities between \system{} and crowd-sourcing systems, \system{}'s use of LLMs provides more degrees of freedom (e.g., the ability to generate code in seconds), creates new sets of challenges (e.g., token optimization), and does not have to deal with ``lazy'' workers, which was one of the key challenges of declarative crowd-based systems. 

